% Part: sets-functions-relations
% Chapter: functions
% Section: basics

\documentclass[../../../include/open-logic-section]{subfiles}

\begin{document}

\olfileid{sfr}{fun}{bas}
\olsection{Dhasar}

\begin{explain}
\emph{Fungsi} yaiku pemetaan kang ngirim saben !!{element}
saka sawijining himpunan menyang !!{element} tartamtu
ing himpunan (liyane) kang wis ditemtokake.
Tuladhane, operasi nambah~$1$ nemtokake fungsi:
saben wilangan~$n$ dipetakake menyang siji-sijine
wilangan~$n+1$.

Kanthi luwih umum, fungsi bisa nampa pasangan,
tupel telu, lan sapanunggalane minangka masukan
lan menehi sawijining asil. Akeh fungsi wis kita
kenal saka aritmetika dhasar. Tuladhane, panambahan
lan ping-pingan iku fungsi. Fungsi kasebut nampa
rong wilangan lan menehi wilangan katelu.

Ing teges matematis lan abstrak iki, fungsi iku
\emph{kothak ireng}: kang wigati mung asil apa
kang dipasangake karo masukan apa, dudu
metodhe kanggo ngetung asile.
\end{explain}

\begin{defn}[Fungsi]
\emph{Fungsi} $f \colon A \to B$ yaiku pemetaan
saben !!{element} saka~$A$ menyang
salah siji !!{element} saka~$B$.

Kita nyebut $A$ minangka \emph{domain} saka~$f$
lan $B$ minangka \emph{kodomain} saka~$f$.
!!{element}s saka~$A$ diarani masukan utawa
\emph{argumen} saka~$f$, lan !!{element}
saka~$B$ kang dipasangake karo argumen~$x$
dening~$f$ diarani \emph{nilai~$f$}
kanggo argumen~$x$, ditulis~$f(x)$.

\emph{Jangkauan} $\ran{f}$ saka~$f$ yaiku
subhimpunan kodomain kang ngemot nilai-nilai~$f$
kanggo sawijining argumen;
$\ran{f} =\Setabs{f(x)}{x \in A}$.
\end{defn}

Diagram ing \olref{fig:function} bisa mbiyantu
ngerteni fungsi. Elips ing sisih kiwa makili
\emph{domain} fungsi; elips ing sisih tengen
makili \emph{kodomain} fungsi; lan panah
nuduhake saka \emph{argumen} ing domain
menyang \emph{nilai} kang cocog ing kodomain.

\begin{figure}
  \olasset{assets/diagrams/function.tikz}
  \caption{Fungsi yaiku pemetaan saben !!{element}
    saka siji himpunan menyang !!a{element}
    saka himpunan liyane. Panah nuduhake
    saka argumen ing domain menyang nilai
    kang cocog ing kodomain.}
  \ollabel{fig:function}
\end{figure}

\begin{ex}
Ping-pingan nampa pasangan wilangan natural
minangka masukan lan metakake menyang wilangan
natural minangka asil, mula saka
$\Nat \times \Nat$ (domain) menyang $\Nat$ (kodomain).
Pranyata jangkauane uga $\Nat$,
amarga saben $n \in \Nat$ iku $n \times 1$.
\end{ex}

\begin{ex}
Ping-pingan iku fungsi amarga masangake saben
masukan---saben pasangan wilangan natural---
karo siji asil: $\times \colon \Nat^2 \to\Nat$.
Dene metakake wilangan natural~$n$
menyang wilangan real~$x$ kang ndadekake
$x^2 = n$ dudu fungsi, amarga saben
wilangan wutuh positif~$n$ duwe rong
oyod kuadrat: $\sqrt{n}$ lan~$-\sqrt{n}$.
Pemetaan iki bisa digawe fungsi kanthi
mulihake oyod kuadrat utama kang ora negatif:
$\sqrt{\phantom{X}} \colon\Nat \to \Real$.
Cathetan koreksi sumber OLFUN-002: tembung Inggris "positive"
ora trep ing nol, amarga domain wilangan natural kalebu nol.
\end{ex}

\begin{ex}
Relasi kang masangake saben siswa ing sawijining
kelas karo biji pungkasane iku fungsi---
ora ana siswa kang bisa oleh rong biji pungkasan
kang beda ing kelas kang padha.
Relasi kang masangake saben siswa ing sawijining
kelas karo wong tuwane dudu fungsi:
siswa bisa ora duwe wong tuwa, duwe loro,
utawa luwih saka loro.
\end{ex}

\begin{explain}
Kita bisa nemtokake fungsi kanthi nyatakake
kanthi tliti apa nilai fungsi kanggo saben
argumen kang bisa ana. Cara-carane bisa
kanthi menehi rumus, njlentrehake metodhe
kanggo ngetung nilai, utawa ndhaptar
nilai kanggo saben argumen.
Kanthi cara apa wae fungsi ditemtokake,
kita kudu mesthekake yen kanggo saben
argumen kita nemtokake siji, lan mung siji, nilai.
\end{explain}


\begin{ex}
Umpamakna $f \colon \Nat \to \Nat$
ditetepake supaya $f(x) = x+1$.
Definisi iki netepake $f$ minangka fungsi
kang nampa wilangan natural lan menehi
wilangan natural minangka asil.
Definisi iki nerangake yen kanggo wilangan
natural~$x$, $f$ menehi paneruse~$x+1$.
Ing kasus iki, kodomain $\Nat$ dudu
jangkauan~$f$, amarga wilangan natural~$0$
dudu panerus saka wilangan natural apa wae.
Jangkauan~$f$ yaiku himpunan kabeh
wilangan wutuh positif, $\Int^{+}$.
\end{ex}

\begin{ex}\ollabel{examplefunext}
Umpamakna $g \colon \Nat \to \Nat$
ditetepake supaya $g(x) = x+2-1$.
Iki nerangake yen $g$ iku fungsi kang
nampa wilangan natural lan menehi
wilangan natural minangka asil.
Kanggo wilangan natural~$x$, $g$ bakal
menehi pandhisik saka panerus saka
paneruse~$x$, yaiku $x+1$.
Cathetan koreksi sumber OLFUN-003: sumber Inggris ngganti jeneng
masukan ing tengah ukara; terjemahan iki nggunakake siji jeneng kanthi ajeg.
\end{ex}

\begin{explain}
Kita lagi wae ngrembug rong fungsi,
$f$ lan $g$, kanthi \emph{definisi}
kang beda. Nanging kalorone iku
\emph{fungsi kang padha}.
Amarga kanggo saben wilangan natural~$n$,
kita duwe $f(n) = n+1 = n+2-1 =g(n)$.
Kanthi tembung liya: definisi kanggo $f$
lan~$g$ netepake pemetaan kang padha
nganggo persamaan kang beda.
Mula kanthi ora langsung kita nggunakake
prinsip ekstensionalitas kanggo fungsi,
\[
  \text{yen }\forall x\, f(x) = g(x)\text{, mula }f = g
\]
angger $f$ lan~$g$ duwe domain lan
kodomain kang padha.
\end{explain}

\begin{ex}
Kita uga bisa nemtokake fungsi miturut
kasus-kasus. Tuladhane, kita bisa
nemtokake $h \colon \Nat \to \Nat$ kanthi
\[
h(x) =
\begin{cases}
  \frac{x}{2} & \text{yen $x$ genep} \\
  \frac{x+1}{2} & \text{yen $x$ ganjil.}
\end{cases}
\]
Amarga saben wilangan natural iku genep
utawa ganjil, asil fungsi iki tansah
wilangan natural.
Elinga yen fungsi ditemtokake miturut kasus,
saben masukan kang bisa ana kudu
kalebu ing persis siji kasus.
Kadhangkala iki mbutuhake bukti yen
kasus-kasus kasebut nyakup kabeh
lan ora ana kang tumpang tindih.
\end{ex}

\end{document}
